% page 271 du fac-simile (image p-270 du scan)
% bloc 172.7 x 254.9 mm ; marge gauche 29.3 mm
\pgc{7.620mm}{0.000mm}{
\l{\cel{55}263}
\l{}
\l{kapsulo. I. (\sou{artilri}o)Alveolo de kupro, qua kontenas materio de-}
\l{\cel{3}toniva, e qua uzesas por komunikar la fairo a la kargajo de}
\l{\cel{3}pulvero di la perkut-armi. - II. (\sou{farmacio}). Envelopilo solv-}
\l{\cel{3}ebla, sen-sapora, aden qua onu inkluzas ula medikamenti por}
\l{\cel{3}disimular lia saporo desagreabla. - DEFIRS.}
\l{}
\l{kaptar.\cel{2}(trans.)\cel{2}I. Sucesar pri sizar ulu,ulo, qua probas eska-}
\l{\cel{3}po. - II. Sizar (la fluo di la liquido de fonto, por evitar la}
\l{\cel{3}infiltri,la alteri, e c.) - EFIS.}
\l{}
\l{kaptacar. (\sou{trans.})\cel{3}Obtenar per manovri ruzoza (donaco, legaco,}
\l{\cel{3}e c.). FS.}
\l{}
\l{kapuco. I. Parto di mantelo, di froko, e c.,forme di boneto tre}
\l{\cel{3}ampla, qua esas charniratre adduktebla adsur la dorso. - II.}
\l{\cel{3}Mantelo kapucoza, quan +weris la viri e la mulieri. - DFIRS.}
\l{}
\l{kapuchino.\cel{2}(\sou{bot.)}\cel{3}Planto klimera, di qua la floro esas oranjo-}
\l{\cel{3}flava, e qua havas, ye sua extremaji suba, apendico kapuco-forma.}
\l{\cel{3}- L.\cel{1}\sou{tropaeolum maius}.}
\l{}
\l{kapucino.(kristanismo). Reguliero ek un de la ordeni franciskana.}
\l{\cel{3}- DEFIRS.}
\l{}
\l{kara.-\cel{2}I. (\sou{pri persono}). Qua inspiras multa tenereso (ad ulu).}
\l{\cel{3}(anke familiare, ma kun senco tre febligita). II. (pri\cel{1}\sou{kozo}).}
\l{\cel{3}Tre prizata pro olua charmo, pro olua importo. - FIS.}
\l{}
\l{karabo. (zool.) Insekto karn-avida, genero ek la serio \textquotedbl{}koleopteri\textquotedbl{},}
\l{\cel{3}e di qua ula speci lansas liquido fetida. - L.\cel{1}\sou{carabus.}\cel{1}-EFIS.}
\l{}
\l{karabino. Fusilo lejera, di qua la tubo havas (ordinare) surfaco}
\l{\cel{3}interna qua esas triizita helicatre, ed extere surfaco taliita}
\l{\cel{4}facioze. - DEFIRS.}
\l{}
\l{\cel{1}\sou{karafo}.\cel{2}Vazo portebla, ordinare sen ansi, kun la ventro globoza,}
\l{\cel{4}la kolo streta, ek kristalo od ek vitro sen-kolora, di qua la}
\l{\cel{4}stopilo esas ek la sama materio, e qua generale destinesis a}
\l{\cel{4}kontenor od a varsor la repast-aquo. - DEFIS.}
\l{}
\l{\cel{1}\sou{karako}. Quaza kazako apertita, quan +weris olim la rurani, la}
\l{\cel{4}soldati.}
\l{}
\l{\cel{1}\sou{karakalo}.\cel{2}(\sou{zool.)}\cel{3}Sorto di kato sovaja. L.\cel{1}\sou{lynx caracal}.-DEFIRS.}
\l{}
\l{\cel{1}\sou{karakolar}. (\sou{netrans}.) (\sou{aludante kavalo}). I. Movar su turne, ad-}
\l{\cel{4}dextre ed adsinistre. - II. Irar kapricoze, per salti, per}
\l{\cel{4}saluto-flexi. - DEFIS.}
\l{}
\l{\cel{1}\sou{karaktero}. I. Traito partikulara, eso-maniero propra, qua distin-}
\l{\cel{4}gas kozo de altra kozo. - II. La eso-maniero morala, psikala,}
\l{\cel{4}spiritala (kontraste kun :la kompreno-fakultato). DEFIRS.}
\l{}
\l{\cel{1}\sou{karakteristiko}.\cel{2}(matem.) La nombro integra di logaritmo.- DEFIS.}
\l{}
\l{\cel{1}\sou{karambolar}.\cel{2}(\sou{netrans}.)(\sou{biliardo-ludo}). Agar la stroko qua}
\l{\cel{4}konsistas ye shokar du buli per la sua. - DFIRS.}
}
